\chapter*{引言}
\markboth{引言}{引言}
\addcontentsline{toc}{chapter}{引言}
\label{intro}
\SourcePage{9}
\noindent\textbf{初步印象。}Fukaya 范畴是与一个辛流形相联系的代数结构。它编码了 Lagrangian 子流形的信息、任意两个此类子流形的相交方式，以及伪全纯圆盘（或多边形）的信息。把如此丰富的数据集中在一个对象之中，使 Fukaya 范畴引人入胜，却也让人有些望而生畏，从辛拓扑的角度来看尤其如此。在继续之前，我们想对这一点稍作展开。下面是一段面向非专家读者的非正式讨论，目的只是让读者对整体情形有一个大致印象；其中有不少并不严格的陈述，讨论范围也比本书其余部分略宽。

固定一个辛流形 $M^{2n}$。粗略地说，Fukaya 范畴的对象是 Lagrangian 子流形 $L^n\subset M$。此类子流形组成的空间是无限维的（局部模型为 $L$ 上闭 $1$ 形式组成的空间）。不过，Hamilton 同痕的 Lagrangian 子流形在 Fukaya 范畴中给出同构的对象，这就把局部自由度降到了第一 Betti 数 $b_1(L)$。事实上，在某些简单情形中，按 Hamilton 同痕分类的问题已经完全解决：例如曲面上的曲线，此时问题本质上属于拓扑；还有较不显然的例子，如 $S^2\times S^2$ 中的 Lagrangian 球面\bcite{75}，以及少数相关例子。另有一些一般性的现象也有助于这一问题：由于存在文献\bcite{60}意义下的障碍，并非所有 Lagrangian 子流形都会实际出现在 Fukaya 范畴中；有些虽会出现，却是平凡的，即零对象；而 Fukaya 范畴中的同构关系也可能比 Lagrangian 同痕更弱。第一点的一个例子，只需考察 $S^2$ 中的圆周：只有当其补集中的两个半球面积相等时，所需的障碍抵消才会发生；更有意思的例子见\bcite{31}、\bcite{32}。第二种现象发生在实点集 $\RP^n\subset\CP^n$（$n>1$）上，条件是它具有 Spin 结构（$n=4k-1$），并且 Fukaya 范畴采用 $\mathbb Q$ 系数。对于第三点，$\CP^2$ 中某些 Lagrangian 环面是候选例子，但据作者所知，其中尚无任何一个得到严格验证。有了这些现象，对于辛拓扑相对简单的流形，原则上计算整个 Fukaya 范畴就成为可行的目标。上面引用的论文，以及其他一些论文，例如\bcite{132}、\bcite{102}、\bcite{101}、\bcite{59}，都可视为朝这一方向推进的、更广泛的持续研究的一部分。然而，对于任何复杂程度稍高的辛流形（例如 $\CP^3$ 中次数至少为 $4$ 的复超曲面），直接分类其 Fukaya 范畴中的对象看来仍遥不可及。

也许只是我们把目标定得太高了。毕竟，即使在纯代数背景下，例如考虑有限维代数上的模，能够实现完全分类的情形也只占极少数\bcite{41}。
\SourcePage{10}
不过，两者有一个显著区别：对于两个明确给出的模，存在一个显然的算法，可以判定它们是否同构。更一般地，给定有限个模，很容易确定它们之间的态射以及态射的复合映射。相比之下，在 Fukaya 范畴中，态射是 Floer 上链群，复合映射则是这些群上的运算；二者都通过适当的伪全纯映射模空间来定义。在某些情形中，可以用几何方法处理这些空间：当 $M$ 是余切丛，而 Lagrangian 子流形 $L=\operatorname{graph}(df)$ 是截面时，一种“绝热极限”方法把问题化为有限维 Morse 理论中的问题\bcite{58}（这一方法的变体也已在若干相关情形中取得成功）；当 $M$ 是定义在实数域上的代数簇，而我们只考虑其实点集 $L=M_{\R}$ 时，可以运用代数几何工具；最后，当 $M$ 是曲面时，对伪全纯曲线的研究在很大程度上成为组合问题。但在这些情形之外，Floer 理论的计算历来依赖逐个例子处理的方法，无法事先保证成功。本书试图应对这一问题。希望本书与其他进展一道，能够使 Fukaya 范畴成为更易于运用的工具。

\noindent\textbf{研究策略。}鉴于上述困难，显然若想有所进展，就必须作出一些大幅简化。首先，我们只考虑正合辛流形 $M$，即辛形式 $\omega$ 是某个 $1$ 形式 $\theta$ 的外微分的流形。这当然意味着 $M$ 不可能是闭流形；不过，我们施加标准的凸性条件，使有限面积的伪全纯曲线感受不到这一缺憾。仿射代数簇提供了丰富的例子。Lagrangian 子流形 $L\subset M$ 中最简单的一类是正合的那些，即 $\theta|L$ 是某个函数的微分；我们将只处理这一类，事先排除其他所有情形。这一限制带来许多有用的技术性结论：例如，正合 Lagrangian 子流形不会有障碍，而且每一个都给出 Fukaya 范畴中的非零对象，其自同态环就是经典的上同调环（没有量子修正）。此外，正合性更根本的好处，是消除了 Fukaya 范畴对辛面积参数（或 Novikov 环参数）的依赖。镜像对称表明，这种依赖往往由非平凡的超越函数给出，而\bcite{108}、\bcite{57}等具体计算也证实了这一点。在我们的情形中，该参数被设为零（或无穷大，视解释而定）；这使问题从计数几何的角度看不那么有趣，却更便于计算。还应提到，最终可以尝试把仿射簇的 Fukaya 范畴作为研究其射影完备化的 Fukaya 范畴的踏脚石\bcite{129}。

\SourcePage{11}
为了使问题更接近经典同调代数，我们还将施加一个 Calabi–Yau 型条件（确切地说，是 $2c_1(M)=0$），并对 Lagrangian 子流形作出相应假设，从而使所有 Floer 群都具有 $\Z$ 分次。由此得到的设置可概括如下。Fukaya 范畴 $\F(M)$ 的对象称为\emph{正合 Lagrangian 膜}，是三元组 $L^{\br}=(L,\alpha^{\br},P^{\br})$，其中 $L$ 是 $M$ 的闭正合 Lagrangian 子流形，$\alpha^{\br}$ 是 $L$ 上称为\emph{分次}的实值函数，而 $P^{\br}$ 是 $L$ 上的一个 Pin 结构。态射空间是 Floer 上链群 $CF^*(L_0^{\br},L_1^{\br})$；它们是有限维分次向量空间（系数取在某个固定但任意的域 $\K$ 中；Pin 结构实际上只在 $\operatorname{char}(\K)\ne2$ 时才必需）。这些空间带有多重线性的复合映射 $\mu^d$，$d\ge1$，并由此构成一个 $\Ai$ 结构。\znote{3}{这里 $\mu^d$ 并非彼此独立：$d=1$ 是微分，$d=2$ 是同伦意义下的乘法，$d=3$ 度量乘法偏离结合律的程度，更高阶项度量下一层的偏差，全部项合起来满足二次的 $\Ai$ 结合律方程。用操作子（operad）的话说，$\Ai$ 操作子是结合操作子的一个 dg 分辨，$\Ai$ 代数就是其上的代数；Fukaya 情形之所以自然出现这种只在同伦意义下成立的结合律，是因为伪全纯曲线模空间只在边界处退化成低阶模空间。[来源：nLab \textquotedblleft A-infinity-algebra\textquotedblright，https:/\allowbreak /\allowbreak ncatlab.org/\allowbreak nlab/\allowbreak show/\allowbreak A-infinity-algebra（2026-11-15 访问）。]}定义需要作出辅助选择，但所得结果在拟同构意义下与这些选择无关。

我们的第二个主要想法来自 Kontsevich 的工作，即对 Fukaya 范畴作形式扩张。这些都是纯代数构造，可以应用于任意 $\Ai$ 范畴 $\A$。第一步是引入链复形的类似物，称为\emph{扭曲复形}；它们组成另一个包含 $\A$ 的 $\Ai$ 范畴 $\Tw\A$。\znote{3}{扭曲复形是链复形的推广：在 dg 情形，一个扭曲复形由一族除有限个外为零的分次对象 $\{E_i\}$ 连同次数为 $i-j+1$ 的态射 $q_{ij}:E_i\to E_j$ 组成，并满足 $dq_{ij}+\sum_k q_{kj}q_{ik}=0$；$q$ 就是所谓\textquotedblleft 扭\textquotedblright，只保留其中一项时便是普通链复形与它的微分。只留一项时 $q$ 给出映射锥，故 $\Tw\A$ 自动含有锥；而令 $q=0$ 便把 $\A$ 嵌入其中。[来源：nLab \textquotedblleft twisted complex\textquotedblright，https:/\allowbreak /\allowbreak ncatlab.org/\allowbreak nlab/\allowbreak show/\allowbreak twisted+\allowbreak complex（2026-11-15 访问）。]}这一范畴含有态射的映射锥，因此具有我们所谓的三角化的 $\Ai$ 范畴结构；其上同调层次的范畴 $H^0(\Tw\A)$，也称为\emph{导出范畴} $D(\A)$，是经典意义下的三角范畴。第二个有用步骤，是取导出范畴的\emph{分裂闭包}（或 Karoubi 完备化）$D^{\pi}(\A)$，把所有幂等自同态的像作为形式直和项加入进来；在 $\Ai$ 层次上也有相应构造，记为 $\Pi(\Tw\A)$。\znote{3}{Karoubi 完备化（幂等完备化）把所有幂等态射 $e:X\to X$ 形式地分裂，得到对象 $(X,e)$ 以及 $X\to(X,e)\to X$；它是使一切幂等态射都可分裂的最小扩张，对任意具备该性质的范畴满足普遍性。$\Pi(\Tw\A)$ 正是这一构造：$\Tw\A$ 只给出锥，再分裂幂等才得到直和项，于是三角闭包与幂等闭包一并完成。$\Tw\A$ 与 $\Pi(\Tw\A)$ 的差别正在于此。[来源：nLab \textquotedblleft Karoubi envelope\textquotedblright，https:/\allowbreak /\allowbreak ncatlab.org/\allowbreak nlab/\allowbreak show/\allowbreak Karoubi+\allowbreak envelope（2026-11-15 访问）。]}每次扩张都会使范畴更容易描述：这一说法表面上似乎矛盾，但许多读者会很清楚它为什么成立。作为说明，考虑\emph{分裂生成元}的概念：若从某一组对象出发，反复取映射锥并分裂出直和项，可以构造出 $\Pi(\Tw\A)$ 的所有对象，就称这组对象分裂生成 $\Pi(\Tw\A)$。只要 $\B\subset\Pi(\Tw\A)$ 是由这样一组分裂生成元构成的全 $\Ai$ 子范畴，$\Pi(\Tw\B)$ 就与 $\Pi(\Tw\A)$ 拟等价。因此，如果只关心 $\Pi(\Tw\A)$，实际上确定一个这样的 $\B$ 就够了。要记住，分裂生成是一个很弱的性质；例如，它并不意味着 Grothendieck 群 $K_0(D^{\pi}(\A))$ 由这些分裂生成元的类所张成。因此，期待找到有限个 Lagrangian 子流形，分裂生成 $\Pi(\Tw\F(M))$，并非不合理；我们将看到，在有利的情形中确实可以做到这一点。

谈到 Fukaya 范畴时使用“计算”一词，意味着必须以某种组合方式编码其基础几何结构，即辛流形和 Lagrangian 子流形。为此，我们将系统地运用 Picard–Lefschetz 理论。一般而言，这一方法在辛几何中的重要性在很大程度上来自 Donaldson 的工作\bcite{38}；这里我们只使用该理论最初等的部分，即单值变换的辛性质，以及 Lefschetz 顶针和消失圈的 Lagrangian 性质。
\SourcePage{12}
为了利用这些性质，我们将依照 Kontsevich 的一个想法，引入 Lefschetz 纤维化的 Fukaya 范畴；从某种意义上说，它介于总空间的通常 Fukaya 范畴与光滑纤维的 Fukaya 范畴之间。随后，我们将以此建立一套按维数归纳进行计算的机制；这套机制最初以猜想的形式提出于\bcite{126}。

\noindent\textbf{计算 Fukaya 范畴。}为简化陈述，我们只考虑代数簇的情形。具体地，设 $X\subset\CC^N$ 是一个 $(n+1)$ 维仿射代数簇，假设它光滑，并且在无穷远处也光滑。后一条件意味着它的闭包 $\overline X\subset\CP^N$ 仍然光滑，且与无穷远超平面 $\CP^{N-1}=\CP^N\setminus\CC^N$ 横截相交。给 $X$ 配备 Fubini–Study Kähler 形式的限制，它便成为正合辛流形。此外，我们假设典范丛 $K_{\overline X}$ 对某个 $d\in\Z$ 同构于 $\mathcal O_{\overline X}(-d)$，从而 $X$ 是 Calabi–Yau 的。一般的仿射超平面截面 $X^1=X\cap\CC^{N-1}$ 也继承所有这些性质。考虑系数取在某个域 $\K$ 中的 Fukaya 范畴 $\F(X)$ 和 $\F(X^1)$。我们先讨论 $X^1$，因为在这一情形中得到的结果略强一些。

\begingroup\renewcommand{\theoremname}{定理 A}
\begin{theorem*}\phantomsection\label{st:A}
假设 $d\ne2$、$n\ge1$ 且 $\operatorname{char}(\K)\ne2$。那么，$\Pi(\Tw\F(X^1))$ 可以由对 $X$ 应用 Picard–Lefschetz 理论所得的组合数据计算出来。\znote{4}{定理 A 的假设在原书印刷页 4 作 $n\ge1$；紧接其后的解释段（印刷页 5）却说 $n>1$ 排除了最低的非平凡维数，并称 $n=1$ 的情形“只是为方便起见”才排除、可用类似方法处理。可见 $n>1$ 是证明叙述中的便利假设，不是定理本身的限制；第 19 节的归纳证明也按 $n\ge1$ 陈述。译文从原书定理陈述，取 $n\ge1$，不将 $n=1$ 排除在定理之外。[来源：原书 p.~4 定理 A、p.~5 说明与 §19 证明；属原书表述不统一，无外部依据。]}
\end{theorem*}\endgroup

这一陈述需要作些解释。首先，对某个 $\Ai$ 范畴 $\A$，称 $\Pi(\Tw\A)$ 可计算，含义如下：可以明确写出另一个在很强意义下有限的 $\Ai$ 范畴 $\B$——其对象只有有限个，态射组成有限维向量空间，且当阶数 $d\gg0$ 足够大时，复合映射 $\mu^d_{\B}$ 全部为零——使 $\Pi(\Tw\B)$ 与 $\Pi(\Tw\A)$ 成为拟等价的 $\Ai$ 范畴。特别地，这意味着 $D^{\pi}(\B)\cong D^{\pi}(\A)$，作为经典三角范畴同构。注意，若合理选择系数域 $\K$（例如有限域或 $\mathbb Q$），$\B$ 确实只包含有限量的信息。可以将这些信息输入计算机程序，通过明确构造扭曲复形及其幂等自同态，生成 $D^{\pi}(\B)$ 中对象的同构类的完整无限列表，并同时给出态射、复合映射和正合三角。当然，仅凭这一点本身，并不能解决有关 $D^{\pi}(\B)$ 的任何重大结构问题。至于“Picard–Lefschetz 理论所产生的组合数据”的含义，精确解释留待\sr{section:19}{第 19 节}；大致想法是考察从 $X$ 到 $\CC$ 和 $\CC^2$ 的线性映射，并用辫单值变换来编码这些映射的分支数据。同样的过程可以应用于 $X^1$，并反复进行，直到维数降至零。给定定义 $X$ 的明确方程，可以借助消元理论，以算法方式提取这些数据
\SourcePage{13}
（不过，除了最简单的例子之外，计算复杂度通常都会高到难以承受）。

还需要简要考察该定理的限制。条件 $n>1$ 排除了最低的非平凡维数，即 $X^1$ 为 Riemann 曲面的情形。这实际上只是为了方便：Riemann 曲面的情形可用相似、事实上还略简单的方法处理，但需要作一些特别调整，这会使讨论复杂化。条件 $\operatorname{char}(\K)\ne2$ 是我们证明方法中的技术性产物，因为证明使用了二重分歧覆盖和 $\Z/2$ 作用；很可能存在其他方法，可以取消这一限制。另一个要求 $d\ne2$ 则略为本质一些，我们把讨论留到后面（见\stref{笔记}{19.8}）。

如前所述，对于 $\F(X)$ 本身，我们得到的结果略弱一些。这并不令人意外：例如，目前没有已知方法可以判定 $X$ 是否含有闭正合 Lagrangian 子流形，从而判定 $\F(X)$ 及其导出范畴是否非平凡。与此相对，在 $X^1$ 的情形中，可以把消失圈作为对象的自然来源。我们所能得到的最好结果如下。

\begingroup\renewcommand{\theoremname}{定理 B}
\begin{theorem*}\phantomsection\label{st:B}
假设 $\operatorname{char}(\K)\ne2$。那么，由对 $X$ 应用 Picard–Lefschetz 理论所得的组合数据，可以构造一个 $\Ai$ 范畴，其中包含一个与 $\F(X)$ 拟等价的全子范畴。
\end{theorem*}\endgroup

这里的 $\Ai$ 范畴具有 $\Tw\B$ 的形式，其中 $\B$ 在与前面相同的意义下是有限的。事实上，按构造 $\B$ 将是有向的，这带来几个值得注意的结论。首先，$\Tw\B$ 已经是分裂闭的，因此它包含的不仅是 $\F(X)$，而是整个 $\Pi(\Tw\F(X))$。其次，有向 $\Ai$ 范畴与态射空间有限维的 dg 范畴拟同构。由于扭曲复形范畴继承这一性质，便有如下推论。

\begingroup\renewcommand{\corollaryname}{推论 C}
\begin{corollary*}\phantomsection\label{st:C}
在\sr{st:B}{定理 B} 的条件下，$\F(X)$ 与一个态射空间有限维的 dg 范畴拟同构。
\end{corollary*}\endgroup

\sr{st:B}{定理 B} 在实际中的用途，取决于我们对映射 $\F(X)\longrightarrow\Tw\B$ 的几何意义了解得有多充分。对于 $X$ 中最简单的一类 Lagrangian 子流形，即能够表示为某个超平面截面 Lefschetz 束的匹配圈的 Lagrangian 球面，这些问题很容易处理，因而确实可以得到计算 Floer 上同调群的算法（严格地说，出于与\sr{st:A}{定理 A} 相同的原因，这里应排除 $n=1$ 的情形）。对于更一般的 Lagrangian 子流形，至少就目前而言，这一陈述还主要具有理论性质。

上述结果的证明将在\sr{section:19}{第 19 节}给出；证明涉及本书各处引入的大部分材料。作为一个自然的副产物，我们建立了消失圈的 Hurwitz 移动与有向 Fukaya 范畴的变异之间的一般关系。
\SourcePage{14}
事实上，我们将给出这一关系的两个证明（\stref{定理}{17.20}和\stref{定理}{18.24}；文献\bcite{127}中最初设想的证明与这两个证明都不相同，迄今仍未发表）。另一个有趣结论，是存在一个计算 Lefschetz 纤维化总空间中 Lagrangian 子流形的 Floer 上同调的谱序列（\stref{推论}{18.27}）。当然，与任何谱序列一样，这只是把问题分解为可见的部分，即起始项，以及隐藏的高阶信息，即微分。在这里，微分表达了这些 Lagrangian 子流形如何由 Lefschetz 顶针构成；“构成”一词应作代数解释，即 Fukaya 范畴中的 Postnikov 分解。

除了结尾处的一处例外，本书并不讨论例子或应用。所幸已有文献弥补了这一不足（若想获得一些线索，可参看\sr{chapter:III}{第三章}开头）。此外，任何认真研究 Fukaya 范畴的读者，最终都必须研读奠基性的文献\bcite{60}；其中讨论的许多几何与代数问题，远远超出了本书建立的框架。
